% !PRINTED 415
\weberpage{415}
\begin{center}
\S.\,110.

\gesperrt{Die Ikosaëdergleichung.}
\end{center}

Das Resultat dieser Betrachtungen ist zunächst, sofern es die allgemeinen
Gleichungen von höherem als dem 4\textsuperscript{ten} Grade betrifft,
ein negatives. Es gelingt nicht, durch rationale Resolventenbildung die
allgemeine Gleichung 5\textsuperscript{ten} Grades auf eine Gleichung mit
\gesperrt{einem} Parameter zu reduciren. Wollen wir gleichwohl dies Ziel noch nicht
aufgeben, so bleibt nichts übrig, als dass wir die Variablen
$x_0, x_1, \ldots, x_{n-1}$ nicht mehr als völlig unabhängige Variable
auffassen, dass wir zwischen ihnen gewisse Gleichungen bestehen lassen.
Wir haben es freilich dann nicht mehr mit der allgemeinen Gleichung des
betreffenden Grades zu thun, sondern mit einer specielleren, bei der die
Co\"efficienten nicht unabhängig von einander sind, und es stellt sich
die zweite Frage ein, wie und durch welche irrationale Functionen sich
die allgemeine Gleichung auf diese specielle zurückführen lässt. Gelingt
dies durch Gleichungen niedrigeren Grades, z. B. durch Quadratwurzeln, so
wird immerhin eine Reduction des Problems erreicht sein.

Wir nehmen also jetzt an, dass die Variablen $x_0, x_1, \ldots, x_{n-1}$
durch eine beliebige Anzahl von Relationen
\weq{1}{A(x_0, x_1, \ldots, x_{n-1}) = 0, \quad B(x_0, x_1, \ldots, x_{n-1}) = 0, \ldots}
mit einander verknüpft seien, worin die $A, B, \ldots$ rationale
Functionen der $x$ sind, die an Anzahl und gegenseitiger Abhängigkeit so
beschaffen sind, dass die Werthe der $x$ nicht numerisch festgelegt,
sondern noch variabel sind. Ausserdem legen wir eine Permutationsgruppe
$\mathfrak{A}$ der $x_0, x_1, \ldots, x_{n-1}$ zu Grunde, und betrachten
die Functionen der $x$, die die Permutationen von $\mathfrak{A}$
gestatten, und nur diese als rational, so dass $\mathfrak{A}$ die
Galois'sche Gruppe der Gleichung ist, deren Wurzeln die
$x_0, x_1, \ldots, x_{n-1}$ sind. \gesperrt{Die Relationen (1) müssen dann so
beschaffen sein, dass auch sie die Permutationen $\mathfrak{A}$
gestatten.}

Wir fragen unter diesen Voraussetzungen nach der Möglichkeit, zwei
rationale Functionen $u, z$ von $x_0, x_1, \ldots, x_{n-1}$ so zu
bestimmen, dass zwischen ihnen eine Gleichung
\weq{2}{F(u, z) = 0}
%CONT>>
